%% kapa: exam=1 round=20 year=2009 period=1 form=A subject=civil_law subject_ko=민법 seq=4 number=4
%% source: Q-Net 감정평가사 기출문제 (cst00309) artlSeq=5024178 file=extracted/1008/5024178/2009년 제20회 감정평가사 1차시험 기출문제/1교시.hwp
%% conversion: kordoc 4.24.51 -> tools/gen1.py (auto); review=proofread-sonnet; flags=none
\begin{question}{4}
일물일권주의에 관한 설명으로 옳지 \uline{않은} 것은?
\begin{choices}
  \item 일물일권주의란 물건의 일부 또는 다수의 물건 위에 하나의 물권이 성립할 수 없다는 원칙이다.
  \item 물건의 일부나 집단 위에 하나의 물권을 인정하여야 할 거래계의 필요가 있고 적절한 공시방법이 마련된 경우에는 물건의 일부나 물건의 집단에 하나의 물권의 성립을 인정하기도 한다.
  \item 명인방법을 갖춘 수목의 집단은 독립한 부동산으로서 저당권의 목적물이 될 수 있다.
  \item 지상권과 지역권은 토지의 일부에도 설정할 수 있다.
  \item 건물의 일부에도 전세권을 설정할 수 있다.
\end{choices}
\end{question}
\answer{3}
